
\section{约定}
\label{sec:1-7}

\par\smallskip\noindent\textup{(Conv.1)}\quad
哈密顿向量场 $X_H$ 由
$dH=\omega(X_H,\mathord\cdot)$ 定义。

\par\smallskip\noindent\textup{(Conv.2)}\quad
若双线性型 $\omega(\mathord\cdot,J\mathord\cdot)$ 为正，则称殆复结构 $J$
是 $\omega$-驯服的；若这一双线性型定义了一个黎曼度量，则称 $J$ 与
$\omega$ 相容（当 $M$ 非紧时，还要求该度量在无穷远处有界）。

\par\smallskip\noindent\textup{(Conv.3)}\quad
设 $\mathbf{x}\in BC[1]$ 是 bar 复形的一个元素。定义余结合余积\liftIndex{I035}
$\Delta\colon BC[1]\to BC[1]\otimes BC[1]$ 如下：
\begin{equation}
  \Delta(x_1\otimes\cdots\otimes x_n)
  =\sum_{k=0}^{n}
  (x_1\otimes\cdots\otimes x_k)
  \otimes
  (x_{k+1}\otimes\cdots\otimes x_n).
  \tag{1.34}\label{eq:1-34}
\end{equation}
考虑它的 $n-1$ 次迭代
\[
  \Delta^{n-1}\colon BC[1]\longrightarrow
  \underbrace{BC[1]\otimes\cdots\otimes BC[1]}_{n\text{ 次}}.
\]
我们写成
\begin{equation}
  \Delta^{n-1}(\mathbf{x})
  =\sum_c
  \mathbf{x}^{n:1}_c\otimes\cdots\otimes\mathbf{x}^{n:n}_c,
  \tag{1.35}\label{eq:1-35}
\end{equation}
其中 $c$ 遍历一个依赖于 $n$ 和 $\mathbf{x}$ 的指标集。当 $n=2$ 时，我们也写成
\begin{equation}
  \Delta(\mathbf{x})=\sum_c\mathbf{x}'_c\otimes\mathbf{x}''_c.
  \tag{1.36}\label{eq:1-36}
\end{equation}

\par\smallskip\noindent\textup{(Conv.4)}\quad
设 $R$ 是带幺元的交换环。我们采用如下版本的诺维科夫环：
\[
  \Lambda_{0,\mathrm{nov}}(R)
  =\left\{
    \left.
    \sum_{i=0}^{\infty}a_iT^{\lambda_i}e^{n_i}
    \ \right|\ 
    a_i\in R,
    \ \lambda_i\in\mathbb R_{\geq0},
    \ n_i\in\mathbb Z,
    \ \lim_{i\to\infty}\lambda_i=\infty
  \right\}.
\]
这里 $T$ 和 $e$ 是形式生成元，$\deg e=2$，$\deg T=0$。我们称它为万有诺维科夫环\liftIndex{I024}。
定义
\[
  \Lambda_{\mathrm{nov}}(R)
  =\left\{
    \left.
    \sum_{i=0}^{\infty}a_iT^{\lambda_i}e^{n_i}
    \ \right|\ 
    a_i\in R,
    \ \lambda_i\in\mathbb R,
    \ n_i\in\mathbb Z,
    \ \lim_{i\to\infty}\lambda_i=\infty
  \right\}.
\]
这是 $\Lambda_{0,\mathrm{nov}}(R)$ 的一个局部化；当 $R$ 是域时，它是
$\Lambda_{0,\mathrm{nov}}(R)$ 的分式域。用 $\Lambda^+_{0,\mathrm{nov}}(R)$
表示 $\Lambda_{0,\mathrm{nov}}(R)$ 的如下理想：
\[
  \Lambda^+_{0,\mathrm{nov}}(R)
  =\left\{
    \left.
    \sum_{i=0}^{\infty}a_iT^{\lambda_i}e^{n_i}
    \in\Lambda_{0,\mathrm{nov}}(R)
    \ \right|\ 
    \lambda_i>0
  \right\}.
\]
注意
\begin{equation}
  \frac{\Lambda_{0,\mathrm{nov}}(R)}
       {\Lambda^+_{0,\mathrm{nov}}(R)}
  \cong R[e,e^{-1}].
  \tag{1.37}\label{eq:1-37}
\end{equation}
分别以 $\Lambda^{(d)}_{0,\mathrm{nov}}(R)$ 和
$\Lambda^{(d)}_{\mathrm{nov}}(R)$ 表示它们次数为 $d$ 的部分。当 $R=\mathbb Q$ 时，
我们略去 $R$，写作 $\Lambda_{0,\mathrm{nov}}$ 等。

\par\smallskip\noindent\textup{(Conv.5)}\quad
我们还考虑环
\[
  \Lambda_0(R)
  =\left\{
    \left.
    \sum_{i=0}^{\infty}a_iT^{\lambda_i}
    \ \right|\ 
    a_i\in R,
    \ \lambda_i\in\mathbb R_{\geq0},
    \ \lim_{i\to\infty}\lambda_i=\infty
  \right\},
\]
以及
\[
  \Lambda_+(R)=\Lambda_0(R)\cap\Lambda^+_{0,\mathrm{nov}}(R),
\]
并考虑局部化 $\Lambda(R)=\Lambda_0(R)[T^{-1}]$。我们把
$\Lambda_0(\mathbb Q)$ 等简记为 $\Lambda_0$ 等。\translatornote{原书此式右端写作 $\Lambda^+_{0,\mathrm{nov}}$，省略了系数环 $(R)$；但左端是在任意系数环 $R$ 上定义 $\Lambda_+(R)$，故补为 $\Lambda^+_{0,\mathrm{nov}}(R)$。}

\par\smallskip\noindent\textup{(Conv.6)}\quad
在 $\Lambda_{0,\mathrm{nov}}(R)$ 上定义滤过
\begin{equation}
  F^\lambda\Lambda_{0,\mathrm{nov}}(R)
  =T^\lambda\Lambda_{0,\mathrm{nov}}(R).
  \tag{1.38}\label{eq:1-38}
\end{equation}
这里 $\lambda\geq0$。它在 $\Lambda_{\mathrm{nov}}(R)$ 上诱导滤过
$F^\lambda\Lambda_{\mathrm{nov}}(R)$，其中 $\lambda\in\mathbb R$；它也在
$\Lambda^+_{0,\mathrm{nov}}(R)$ 上诱导一个滤过。这些滤过在
$\Lambda_{0,\mathrm{nov}}(R)$、$\Lambda_{\mathrm{nov}}(R)$ 和
$\Lambda^+_{0,\mathrm{nov}}(R)$ 上诱导度量，而这些环关于该度量完备。
我们以类似方式在 $\Lambda_0(R)$、$\Lambda_+(R)$ 和 $\Lambda(R)$ 上定义滤过。
